% Generated from source-md/3-chapitre2.md. Keep content changes in the Markdown source.

\chapter{Analyse et préparation du projet}
\label{chap:analyse-preparation-projet}


\section{Introduction}


Ce chapitre présente la phase d’analyse et de préparation du projet FESTY. Il a pour objectif de définir les besoins de la plateforme, d’identifier les acteurs qui interagissent avec le système et de préparer les éléments nécessaires à la conception et au développement.


Nous commençons par la spécification des besoins, à travers l’identification des acteurs principaux ainsi que la présentation des besoins fonctionnels et non fonctionnels. Ensuite, nous présentons le pilotage du projet avec Scrum, en précisant l’équipe Scrum, le backlog produit et la planification des sprints.


Ce chapitre présente également les diagrammes de cas d’utilisation globaux de la solution. Enfin, nous décrivons l’environnement de travail, les technologies utilisées ainsi que l’architecture générale de la plateforme.


\section{Spécification des besoins}


La spécification des besoins constitue une étape importante dans la préparation du projet. Elle permet de définir les fonctionnalités attendues de la plateforme ainsi que les contraintes que le système doit respecter.


Dans le cadre du projet FESTY, les besoins sont organisés autour des différents acteurs qui interagissent avec la plateforme. Cette section présente les acteurs principaux, les besoins fonctionnels et les besoins non fonctionnels du système.


\subsection{Identification des acteurs}


Un acteur représente une entité externe qui interagit avec le système afin d’atteindre un objectif précis. Dans le cadre de FESTY, les acteurs correspondent aux différents profils qui utilisent la plateforme selon des rôles et des droits distincts.


Nous distinguons cinq acteurs principaux : le Visiteur, l’Utilisateur, le Partenaire, l’Agent de scan et l’Administrateur.


\begin{table}[H]
\centering
\caption{Identification des acteurs principaux de FESTY}
\label{tableau:2-1}
\renewcommand{\arraystretch}{1.2}
\begin{tabularx}{\textwidth}{|>{\raggedright\arraybackslash}X|>{\raggedright\arraybackslash}X|}
\hline
\rowcolor[HTML]{ECF4FF}
\textbf{Acteur} & \textbf{Description} \\
\hline
Visiteur & Représente une personne non authentifiée qui accède aux fonctionnalités publiques de la plateforme. Il peut explorer les événements et les artistes, consulter les informations disponibles et s’inscrire à la plateforme. \\
\hline
Utilisateur & Représente une personne authentifiée disposant d’un compte sur la plateforme. Il peut gérer son profil et ses préférences, gérer ses billets, interagir autour des événements et soumettre une demande pour devenir partenaire organisateur. \\
\hline
Partenaire & Représente un organisateur d’événements. Il dispose d’un espace dédié et d’un portail partenaire lui permettant de gérer son profil professionnel, ses événements, sa billetterie, ses sessions de scan, son activité et sa monétisation. \\
\hline
Agent de scan & Représente un acteur chargé de contrôler l’accès à un événement. Il peut vérifier les billets des participants à travers le scan et valider leur entrée. \\
\hline
Administrateur & Représente le responsable de la supervision de la plateforme. Il assure la gestion des utilisateurs et des partenaires, le contrôle de l’offre événementielle, la gestion du catalogue artistique, le traitement des opérations financières ainsi que la supervision des risques et des audits. \\
\hline
\end{tabularx}
\end{table}


Dans le cadre de FESTY, le terme « artiste » désigne les stars, chanteurs, DJ ou intervenants associés aux événements.


\subsection{Besoins fonctionnels}


Les besoins fonctionnels décrivent les services que la plateforme doit offrir à ses différents acteurs. Ils permettent de préciser les principales fonctionnalités attendues du système.


Le tableau suivant présente les besoins fonctionnels principaux de FESTY selon chaque acteur.


\begin{table}[H]
\centering
\caption{Besoins fonctionnels principaux de FESTY}
\label{tableau:2-2}
\renewcommand{\arraystretch}{1.2}
\begin{tabularx}{\textwidth}{|>{\raggedright\arraybackslash}X|>{\raggedright\arraybackslash}X|}
\hline
\rowcolor[HTML]{ECF4FF}
\textbf{Acteur} & \textbf{Besoins fonctionnels principaux} \\
\hline
Visiteur & Explorer les événements et les artistes, consulter les informations publiques, s’inscrire à la plateforme et s’authentifier. \\
\hline
Utilisateur & Gérer son profil et ses préférences, interagir autour des événements, gérer sa billetterie et soumettre une demande pour devenir partenaire organisateur. \\
\hline
Partenaire & Gérer son profil partenaire, gérer l’exploitation de ses événements, gérer le line-up, suivre son activité, consulter sa monétisation et organiser le contrôle d’accès aux événements. \\
\hline
Agent de scan & Contrôler l’accès à un événement en vérifiant les billets des participants et en validant leur entrée. \\
\hline
Administrateur & Gérer les utilisateurs et les partenaires, contrôler l’offre événementielle, gérer le catalogue artistique, traiter les opérations financières et superviser les risques et les audits. \\
\hline
\end{tabularx}
\end{table}


Ces besoins sont présentés à un niveau global afin de donner une vue synthétique du périmètre fonctionnel de la plateforme. Les détails de chaque fonctionnalité seront précisés dans les chapitres de sprint à travers les diagrammes de cas d’utilisation détaillés, les descriptions textuelles et les éléments de conception.


\subsection{Besoins non fonctionnels}


Les besoins non fonctionnels décrivent les contraintes de qualité que doit respecter la plateforme. Ils ne représentent pas directement des fonctionnalités visibles par l’utilisateur, mais ils sont essentiels pour garantir la fiabilité, la sécurité et la bonne utilisation du système.


\begin{table}[H]
\centering
\caption{Besoins non fonctionnels de FESTY}
\label{tableau:2-3}
\renewcommand{\arraystretch}{1.2}
\begin{tabularx}{\textwidth}{|>{\raggedright\arraybackslash}X|>{\raggedright\arraybackslash}X|}
\hline
\rowcolor[HTML]{ECF4FF}
\textbf{Besoin non fonctionnel} & \textbf{Description} \\
\hline
Sécurité & La plateforme doit protéger les comptes utilisateurs, les données personnelles, les billets numériques, les paiements et les accès aux fonctionnalités sensibles. \\
\hline
Performance & Le système doit assurer un temps de réponse acceptable lors de la consultation des événements, de la réservation et du paiement des billets ainsi que du scan à l’entrée des événements. \\
\hline
Disponibilité & Les fonctionnalités principales doivent rester accessibles afin de garantir une utilisation continue par les visiteurs, utilisateurs, partenaires et agents de scan. \\
\hline
Ergonomie & La plateforme doit proposer des interfaces claires, intuitives et adaptées à l’application mobile, au back-office et au portail partenaire. \\
\hline
Compatibilité & La solution doit être compatible avec les environnements mobiles et web ciblés. \\
\hline
Maintenabilité & L’architecture doit faciliter l’évolution de la plateforme, l’ajout de nouvelles fonctionnalités et la correction des anomalies. \\
\hline
Scalabilité & Le système doit pouvoir supporter l’augmentation du nombre d’utilisateurs, d’événements, de billets vendus et d’opérations de scan. \\
\hline
Traçabilité & Les opérations sensibles, telles que les paiements, les validations, les scans et les actions administratives, doivent être suivies et historisées. \\
\hline
Confidentialité & Les informations personnelles des utilisateurs, les données des partenaires et les informations liées aux transactions doivent être protégées contre tout accès non autorisé. \\
\hline
\end{tabularx}
\end{table}


Ces contraintes permettent d’orienter les choix techniques et architecturaux de la solution. Elles seront prises en compte lors de la conception, du développement et de la validation des différents modules de la plateforme.


\section{Pilotage du projet avec Scrum}


Le pilotage du projet avec Scrum permet d’organiser le développement de FESTY de manière progressive et itérative. Après le choix de Scrum présenté dans le chapitre précédent, cette section précise son application dans le cadre du projet.


Dans notre cas, Scrum est utilisé afin de découper le projet en plusieurs sprints cohérents. Chaque sprint regroupe un ensemble de fonctionnalités liées à un objectif précis, ce qui permet de suivre l’avancement du développement, de prioriser les besoins et de construire progressivement la plateforme.


\subsection{Équipe et rôles Scrum}


Dans le cadre du projet FESTY, l’organisation Scrum a été adaptée au contexte d’un projet de fin d’études. Les rôles ont été définis selon les responsabilités liées au suivi, à l’encadrement et à la réalisation du projet.


\begin{table}[H]
\centering
\caption{Équipe et rôles du projet FESTY}
\label{tableau:2-4}
\renewcommand{\arraystretch}{1.2}
\begin{tabularx}{\textwidth}{|>{\raggedright\arraybackslash}X|>{\raggedright\arraybackslash}X|>{\raggedright\arraybackslash}X|}
\hline
\rowcolor[HTML]{ECF4FF}
\textbf{Rôle dans le projet} & \textbf{Responsable} & \textbf{Missions} \\
\hline
Product Owner & Encadrant professionnel / représentant du besoin métier & Définir les besoins fonctionnels, valider les priorités et orienter les choix liés aux fonctionnalités de la plateforme. \\
\hline
Scrum Master & Étudiant réalisateur du projet & Organiser le travail, suivre l’avancement, identifier les difficultés et assurer l’application de la démarche Scrum. \\
\hline
Équipe de développement & Étudiant réalisateur du projet & Concevoir, développer, tester et intégrer les différentes fonctionnalités de la plateforme FESTY. \\
\hline
Encadrant académique & Encadrant universitaire & Assurer le suivi académique, valider la cohérence méthodologique du rapport et orienter la structuration du travail. \\
\hline
\end{tabularx}
\end{table}


Cette répartition permet d’adapter Scrum au contexte réel du projet. L’étudiant assure à la fois le rôle de Scrum Master et de membre de l’équipe de développement, tandis que les encadrants assurent le suivi, l’orientation et la validation du travail réalisé.


\subsection{Backlog du produit}


Le backlog du produit regroupe l’ensemble des fonctionnalités à réaliser dans le cadre du projet FESTY. Il permet de structurer les besoins sous forme de modules et de user stories, tout en définissant leur niveau de priorité.


Afin de prioriser les fonctionnalités, nous avons adopté la méthode MoSCoW [5]. Cette méthode permet de classer les besoins selon leur importance dans la version actuelle du produit.


\begin{table}[H]
\centering
\caption{Signification des priorités MoSCoW}
\label{tableau:2-5}
\renewcommand{\arraystretch}{1.2}
\begin{tabularx}{\textwidth}{|>{\raggedright\arraybackslash}X|>{\raggedright\arraybackslash}X|}
\hline
\rowcolor[HTML]{ECF4FF}
\textbf{Priorité} & \textbf{Signification} \\
\hline
\textbf{M — Must Have} & Fonctionnalité indispensable au bon fonctionnement de la plateforme. \\
\hline
\textbf{S — Should Have} & Fonctionnalité importante, mais pouvant être livrée après les fonctionnalités indispensables. \\
\hline
\textbf{C — Could Have} & Fonctionnalité utile ou complémentaire, mais non prioritaire dans la version actuelle. \\
\hline
\textbf{W — Won’t Have} & Fonctionnalité exclue du périmètre actuel, pouvant être envisagée dans une version future. \\
\hline
\end{tabularx}
\end{table}


Le tableau suivant présente le Product Backlog priorisé de la plateforme FESTY.


{\footnotesize
\renewcommand{\arraystretch}{1.2}
\setlength{\tabcolsep}{4pt}
\begin{xltabular}{\textwidth}{|>{\centering\arraybackslash}p{0.8cm}|>{\raggedright\arraybackslash}p{2.7cm}|>{\centering\arraybackslash}p{1.2cm}|>{\raggedright\arraybackslash}X|>{\centering\arraybackslash}p{1.5cm}|}
\caption{Backlog du produit priorisé selon la méthode MoSCoW}
\label{tableau:2-6}\\
\hline
\rowcolor[HTML]{ECF4FF}
\textbf{ID} & \textbf{Module} & \textbf{ID Story} & \textbf{User Story} & \textbf{Priorité} \\
\hline
\endfirsthead

\hline
\rowcolor[HTML]{ECF4FF}
\textbf{ID} & \textbf{Module} & \textbf{ID Story} & \textbf{User Story} & \textbf{Priorité} \\
\hline
\endhead

\hline
\multicolumn{5}{r}{\textit{Suite de la page suivante}} \\
\hline
\endfoot

\hline
\endlastfoot

1 & Authentification et profil & 1.1 & En tant que visiteur, je veux créer un compte afin d’accéder à la plateforme. & M \\
\hline
1 & Authentification et profil & 1.2 & En tant qu’utilisateur, je veux m’authentifier afin d’accéder à mon espace personnel. & M \\
\hline
1 & Authentification et profil & 1.3 & En tant qu’utilisateur, je veux réinitialiser mon mot de passe afin de récupérer l’accès à mon compte. & S \\
\hline
1 & Authentification et profil & 1.4 & En tant qu’administrateur, je veux m’authentifier au back-office afin d’accéder aux fonctionnalités d’administration de la plateforme. & M \\
\hline
1 & Authentification et profil & 1.5 & En tant qu’utilisateur, je veux consulter et modifier mon profil afin de maintenir mes informations à jour. & M \\
\hline
2 & Partenaires et événements & 2.1 & En tant qu’utilisateur, je veux soumettre une demande pour devenir partenaire organisateur. & M \\
\hline
2 & Partenaires et événements & 2.2 & En tant qu’administrateur, je veux traiter les demandes partenaires afin de valider ou refuser l’accès à l’espace partenaire. & M \\
\hline
2 & Partenaires et événements & 2.3 & En tant que partenaire, je veux gérer mon profil partenaire afin de présenter mes informations professionnelles. & S \\
\hline
2 & Partenaires et événements & 2.4 & En tant que partenaire, je veux créer et gérer mes événements afin de maintenir leurs informations à jour. & M \\
\hline
2 & Partenaires et événements & 2.5 & En tant qu’administrateur, je veux contrôler l’offre événementielle afin de garantir la conformité du contenu. & M \\
\hline
3 & Exploration, personnalisation, artistes et interactions & 3.1 & En tant que visiteur ou utilisateur, je veux explorer les événements afin de découvrir les offres disponibles. & M \\
\hline
3 & Exploration, personnalisation, artistes et interactions & 3.2 & En tant que visiteur ou utilisateur, je veux consulter les artistes associés et le line-up des événements afin d’obtenir plus d’informations. & S \\
\hline
3 & Exploration, personnalisation, artistes et interactions & 3.3 & En tant que partenaire, je veux gérer le line-up de mes événements afin d’associer des artistes aux événements publiés. & M \\
\hline
3 & Exploration, personnalisation, artistes et interactions & 3.4 & En tant qu’utilisateur, je veux gérer mes préférences afin de personnaliser mon expérience. & M \\
\hline
3 & Exploration, personnalisation, artistes et interactions & 3.5 & En tant qu’utilisateur, je veux interagir autour des événements afin d’enrichir mon expérience. & M \\
\hline
3 & Exploration, personnalisation, artistes et interactions & 3.6 & En tant qu’administrateur, je veux gérer le catalogue artistique afin de contrôler les stars disponibles et les propositions soumises par les partenaires. & M \\
\hline
4 & Billetterie, paiement, scan et supervision & 4.1 & En tant qu’utilisateur, je veux réserver et payer un billet afin de participer à un événement. & M \\
\hline
4 & Billetterie, paiement, scan et supervision & 4.2 & En tant qu’utilisateur, je veux consulter mes tickets afin d’accéder aux informations de mes achats. & M \\
\hline
4 & Billetterie, paiement, scan et supervision & 4.3 & En tant qu’utilisateur, je veux demander un remboursement lorsque cela est possible. & S \\
\hline
4 & Billetterie, paiement, scan et supervision & 4.4 & En tant qu’utilisateur, je veux mettre un ticket en revente ou acheter un ticket revendu via la marketplace. & S \\
\hline
4 & Billetterie, paiement, scan et supervision & 4.5 & En tant qu’utilisateur, je veux proposer un ticket externe à la vente afin de le publier dans la marketplace. & C \\
\hline
4 & Billetterie, paiement, scan et supervision & 4.6 & En tant que partenaire, je veux gérer les sessions de scan et affecter des agents afin d’organiser le contrôle d’accès. & M \\
\hline
4 & Billetterie, paiement, scan et supervision & 4.7 & En tant qu’agent de scan, je veux vérifier les tickets et valider l’entrée d’un participant afin de contrôler l’accès à un événement. & M \\
\hline
4 & Billetterie, paiement, scan et supervision & 4.8 & En tant que partenaire, je veux suivre mon activité et ma monétisation afin de consulter les performances de mes événements. & S \\
\hline
4 & Billetterie, paiement, scan et supervision & 4.9 & En tant qu’administrateur, je veux traiter les opérations financières et superviser les opérations sensibles afin d’assurer le bon fonctionnement de la plateforme. & S \\
\hline
\end{xltabular}
}


Les éléments du Product Backlog représentent les fonctionnalités principales de FESTY. Ils seront ensuite utilisés pour organiser la planification des sprints selon les priorités et l’ordre de réalisation retenu. Ils seront également représentés de manière synthétique dans les diagrammes de cas d’utilisation globaux.


Le Product Backlog présente une vue synthétique des fonctionnalités principales de la plateforme. Les Sprint Backlogs détaillés, présentés dans les chapitres suivants, peuvent décomposer ces fonctionnalités en user stories plus précises selon le périmètre de chaque sprint.


Dans le cadre de cette version, aucune fonctionnalité n’a été classée comme \textbf{Won’t Have}, car le backlog présenté regroupe uniquement les fonctionnalités retenues dans le périmètre du projet.


\subsection{Planification des sprints}


La planification des sprints permet d’organiser le développement de la plateforme FESTY en quatre itérations cohérentes. Chaque sprint regroupe un ensemble de fonctionnalités majeures issues du Product Backlog.


Le tableau suivant présente une vue synthétique de cette planification.


\begin{table}[H]
\centering
\caption{Planification des sprints du projet FESTY}
\label{tableau:2-7}
\renewcommand{\arraystretch}{1.2}
\begin{tabularx}{\textwidth}{|>{\raggedright\arraybackslash}X|>{\raggedright\arraybackslash}X|>{\raggedright\arraybackslash}X|>{\raggedright\arraybackslash}X|}
\hline
\rowcolor[HTML]{ECF4FF}
\textbf{Sprint} & \textbf{Objectif} & \textbf{Acteurs principaux} & \textbf{Fonctionnalités principales} \\
\hline
Sprint 1 & Authentification et gestion des accès & Visiteur, Utilisateur, Administrateur & Inscription, vérification OTP, authentification à la plateforme, authentification au back-office avec 2FA et gestion du profil utilisateur. \\
\hline
Sprint 2 & Gestion des partenaires et des événements & Utilisateur, Partenaire, Administrateur & Soumission et traitement des demandes partenaires, gestion du profil partenaire, création et gestion des événements, contrôle de l’offre événementielle. \\
\hline
Sprint 3 & Exploration, personnalisation et interactions autour des événements & Visiteur, Utilisateur, Partenaire, Administrateur & Exploration des événements et des stars, consultation des artistes et du line-up, gestion des préférences, recommandations, favoris, suivi des stars, chat événementiel, gestion du line-up artistique et gestion du catalogue artistique. \\
\hline
Sprint 4 & Billetterie, paiement, revente et contrôle d’accès & Utilisateur, Partenaire, Agent de scan, Administrateur & Réservation des billets, paiement, consultation et gestion des tickets, remboursement, revente, marketplace, ticket externe, scan et validation des accès, suivi partenaire, monétisation et supervision des opérations sensibles. \\
\hline
\end{tabularx}
\end{table}


Cette planification donne une vue globale de l’organisation du développement. Les détails de chaque sprint, notamment le Sprint Backlog, les diagrammes de cas d’utilisation détaillés, les descriptions textuelles, les diagrammes de séquence et les éléments de réalisation, seront présentés dans les chapitres consacrés aux sprints.


\section{Diagrammes de cas d’utilisation globaux}


Afin de présenter une vue générale de la solution, nous distinguons deux diagrammes de cas d’utilisation globaux : l’un pour la plateforme FESTY et l’autre pour le back-office administrateur.


\subsection{Diagramme de cas d’utilisation global de la plateforme FESTY}


Ce diagramme présente les principales fonctionnalités accessibles aux acteurs métier de la plateforme FESTY. Il conserve les relations include vers l’authentification afin de garder une lecture cohérente des accès à la plateforme.


\begin{figure}[H]
\centering
\includegraphics[width=0.95\textwidth]{\detokenize{images/chapitre2/Figure 2.1 Diagramme de cas d’utilisation global de la plateforme FESTY.png}}
\caption{Diagramme de cas d’utilisation global de la plateforme FESTY.}
\label{figure:2-1}
\end{figure}


La figure 2.1 montre que le Visiteur peut explorer les événements, s’inscrire à la plateforme et accéder à l’authentification. Après connexion, l’Utilisateur peut gérer son profil, ses préférences, sa billetterie et ses interactions. Le Partenaire accède à des fonctionnalités dédiées à la gestion de son profil partenaire, à l’exploitation de ses événements, au suivi de son activité et de sa monétisation. L’Agent de scan assure le contrôle d’accès aux événements à travers la vérification des billets des participants.


\subsection{Diagramme de cas d’utilisation global du back-office administrateur}


Ce diagramme présente les principales fonctionnalités du back-office administrateur de FESTY. Il regroupe les opérations de supervision, de gestion et de contrôle accessibles après authentification au back-office.


\begin{figure}[H]
\centering
\includegraphics[width=0.95\textwidth]{\detokenize{images/chapitre2/Figure 2.2 Diagramme de cas d’utilisation global du back-office administrateur.png}}
\caption{Diagramme de cas d’utilisation global du back-office administrateur.}
\label{figure:2-2}
\end{figure}


La figure 2.2 présente les fonctions de gestion des utilisateurs et des partenaires, le contrôle de l’offre événementielle, la gestion du catalogue artistique, le traitement des opérations financières ainsi que la supervision des risques et des audits.


\section{Environnement de travail}


L’environnement de travail regroupe les ressources matérielles et logicielles utilisées pour concevoir, développer, tester et documenter la plateforme FESTY. Cette section présente les principaux éléments ayant permis d’assurer le bon déroulement du projet.


\subsection{Environnement matériel}


Le développement de FESTY a été réalisé à l’aide d’un ordinateur portable utilisé pour le développement, les tests et la modélisation. Des appareils mobiles ont également été utilisés pour vérifier le comportement de l’application dans un contexte proche de l’utilisation réelle.


\begin{table}[H]
\centering
\caption{Environnement matériel utilisé}
\label{tableau:2-8}
\renewcommand{\arraystretch}{1.2}
\begin{tabularx}{\textwidth}{|>{\raggedright\arraybackslash}X|>{\raggedright\arraybackslash}X|}
\hline
\rowcolor[HTML]{ECF4FF}
\textbf{Caractéristique} & \textbf{Description} \\
\hline
Nom de l’appareil & DELL G15 5510 \\
\hline
Processeur & Intel(R) Core(TM) i5-10500H CPU @ 2.50GHz \\
\hline
Mémoire RAM & 24 Go \\
\hline
Carte graphique dédiée & NVIDIA GeForce GTX 1650 — 4 Go \\
\hline
Carte graphique intégrée & Intel(R) UHD Graphics \\
\hline
Stockage & SSD 238 Go \\
\hline
Système d’exploitation & Windows \\
\hline
Appareil de test mobile & Téléphone Android réel ; tests iOS réalisés via TestFlight \\
\hline
\end{tabularx}
\end{table}


\subsection{Environnement logiciel}


Le tableau suivant présente les principaux outils logiciels utilisés durant la réalisation du projet. Ces outils couvrent le développement, les tests, la modélisation, la gestion du code source et la documentation.


\begin{table}[H]
\centering
\caption{Environnement logiciel utilisé}
\label{tableau:2-9}
\renewcommand{\arraystretch}{1.2}
\begin{tabularx}{\textwidth}{|>{\raggedright\arraybackslash}X|>{\raggedright\arraybackslash}X|}
\hline
\rowcolor[HTML]{ECF4FF}
\textbf{Outil} & \textbf{Utilisation} \\
\hline
Visual Studio Code & Édition du code source, des fichiers de configuration et de la documentation du projet. \\
\hline
Android Studio & Compilation, exécution et test de l’application mobile Android. \\
\hline
Postman & Test et validation des API REST. \\
\hline
Git & Gestion locale des versions du code source. \\
\hline
GitHub / GitLab & Hébergement du code source et suivi des modifications. \\
\hline
pgAdmin & Administration et consultation de la base de données PostgreSQL. \\
\hline
Docker & Exécution de certains services nécessaires aux tests et à l’environnement de développement. \\
\hline
Stripe CLI & Test local des webhooks Stripe et simulation des événements de paiement. \\
\hline
Draw.io / PlantUML & Création et amélioration visuelle des diagrammes. \\
\hline
\end{tabularx}
\end{table}


\section{Technologies utilisées}


Cette section présente les principales technologies utilisées pour la réalisation de la plateforme FESTY. Elles sont regroupées selon leur rôle dans le projet : technologies Front-End et Mobile, technologies Back-End, base de données et services externes.


\subsection{Technologies Front-End et Mobile}


La partie Front-End de la plateforme FESTY comprend l’application mobile ainsi que les interfaces web du back-office administrateur et du portail partenaire.


\begin{table}[H]
\centering
\caption{Technologies Front-End et Mobile utilisées}
\label{tableau:2-10}
\renewcommand{\arraystretch}{1.2}
\begin{tabularx}{\textwidth}{|>{\raggedright\arraybackslash}X|>{\raggedright\arraybackslash}X|}
\hline
\rowcolor[HTML]{ECF4FF}
\textbf{Technologie} & \textbf{Rôle dans le projet} \\
\hline
Flutter & Framework utilisé pour développer l’application mobile multiplateforme de FESTY. \\
\hline
Dart [8] & Langage utilisé avec Flutter pour développer les interfaces mobiles et la logique côté client. \\
\hline
React.js [9] & Bibliothèque utilisée pour développer les interfaces web, notamment le tableau de bord administrateur et le portail partenaire. \\
\hline
TypeScript [10] / JavaScript [11] & Langages utilisés pour structurer la logique des interfaces web et améliorer la maintenabilité du code Front-End. \\
\hline
\end{tabularx}
\end{table}


\subsection{Technologies Back-End}


Le Back-End de FESTY est développé avec Spring Boot et organisé selon une architecture modulaire. La base de données utilisée est PostgreSQL. Les services externes tels que Stripe, Firebase Cloud Messaging, Twilio et Bunny.net interviennent dans des fonctionnalités spécifiques.


Cette partie assure la gestion de la logique métier, la sécurité, l’exposition des API REST, la persistance des données et l’intégration avec les services externes. Elle repose sur une architecture monolithique modulaire afin de séparer les domaines fonctionnels tout en conservant une plateforme centralisée.


\begin{table}[H]
\centering
\caption{Technologies Back-End utilisées}
\label{tableau:2-11}
\renewcommand{\arraystretch}{1.2}
\begin{tabularx}{\textwidth}{|>{\raggedright\arraybackslash}X|>{\raggedright\arraybackslash}X|}
\hline
\rowcolor[HTML]{ECF4FF}
\textbf{Technologie} & \textbf{Rôle dans le projet} \\
\hline
Java 17 [12] & Langage principal utilisé pour le développement de la partie Back-End. \\
\hline
Spring Boot [13] & Framework principal utilisé pour développer les services Back-End et exposer les API REST. \\
\hline
Spring Security / JWT & Technologies utilisées pour gérer l’authentification, la sécurité et les autorisations selon les rôles. \\
\hline
Spring Data JPA / Hibernate & Technologies utilisées pour assurer la persistance des données et la communication avec la base PostgreSQL. \\
\hline
Spring Modulith [18] & Technologie utilisée pour structurer le Back-End en modules métier cohérents. \\
\hline
Flyway [19] & Outil utilisé pour gérer les migrations SQL et suivre l’évolution du schéma de la base de données. \\
\hline
Springdoc OpenAPI [20] / Swagger UI & Outil utilisé pour documenter et tester les API REST exposées par le Back-End. \\
\hline
\end{tabularx}
\end{table}


Le choix de Spring Boot permet de développer un Back-End structuré et évolutif. Spring Security et JWT assurent l’authentification interne, la gestion des jetons d’accès et de rafraîchissement ainsi que le contrôle des autorisations selon les rôles, tandis que Spring Data JPA [16] et Hibernate [17] facilitent la gestion des données. Spring Modulith permet d’organiser le code en modules métier, notamment l’authentification, les événements, la billetterie, le paiement, le scan et l’administration.


\subsection{Base de données}


La plateforme FESTY utilise une base de données relationnelle afin de stocker les informations nécessaires au fonctionnement du système, notamment les utilisateurs, les partenaires, les événements, les billets, les paiements, les opérations de scan et les données d’administration.


\begin{table}[H]
\centering
\caption{Technologies liées à la base de données}
\label{tableau:2-12}
\renewcommand{\arraystretch}{1.2}
\begin{tabularx}{\textwidth}{|>{\raggedright\arraybackslash}X|>{\raggedright\arraybackslash}X|}
\hline
\rowcolor[HTML]{ECF4FF}
\textbf{Technologie} & \textbf{Rôle dans le projet} \\
\hline
PostgreSQL & Système de gestion de base de données relationnelle utilisé pour stocker les données principales de la plateforme. \\
\hline
Supabase & Plateforme utilisée pour l’hébergement de la base PostgreSQL. \\
\hline
PostGIS & Extension de PostgreSQL utilisée pour gérer les données géographiques liées aux événements. \\
\hline
Hibernate Spatial [24] & Extension utilisée avec Hibernate pour manipuler les données géographiques côté Back-End. \\
\hline
\end{tabularx}
\end{table}


Dans cette architecture, Supabase [22] est utilisé uniquement comme environnement d’hébergement de la base PostgreSQL. L’authentification, la gestion des utilisateurs, les rôles et les sessions sécurisées sont assurés par le Back-End FESTY à travers Spring Security et JWT.


\subsection{Services externes}


FESTY intègre plusieurs services externes afin de prendre en charge certaines fonctionnalités spécifiques, notamment les paiements, les notifications, les emails, la vérification téléphonique et la gestion des fichiers.


\begin{table}[H]
\centering
\caption{Services externes intégrés à FESTY}
\label{tableau:2-13}
\renewcommand{\arraystretch}{1.2}
\begin{tabularx}{\textwidth}{|>{\raggedright\arraybackslash}X|>{\raggedright\arraybackslash}X|}
\hline
\rowcolor[HTML]{ECF4FF}
\textbf{Service externe} & \textbf{Rôle dans le projet} \\
\hline
Stripe [25] & Gestion des paiements en ligne, des remboursements, des webhooks, des payouts et des comptes connectés des partenaires. \\
\hline
Firebase Cloud Messaging & Envoi des notifications push aux utilisateurs de l’application mobile. \\
\hline
Brevo [27] & Envoi des SMS pour la vérification OTP et des emails transactionnels liés aux comptes, aux notifications et aux opérations importantes. \\
\hline
Bunny.net & Stockage et diffusion des fichiers ou médias utilisés par la plateforme. \\
\hline
\end{tabularx}
\end{table}


\section{Architecture de la solution}


L’architecture de la solution permet de représenter l’organisation générale de la plateforme FESTY. Elle décrit les principaux composants du système, leurs responsabilités ainsi que les échanges entre les différentes couches de la plateforme.


Dans le cadre de ce projet, nous avons adopté une architecture organisée autour d’une couche présentation, d’une couche métier, d’une couche de données et de services externes. Cette séparation permet de faciliter la maintenance, l’évolution du système et la répartition des responsabilités entre les composants.


\subsection{Architecture physique}


L’architecture physique décrit les composants déployés ainsi que les environnements dans lesquels ils s’exécutent. Elle met en évidence les clients de la plateforme, le Back-End, la base de données et les services externes intégrés à la solution.


La plateforme FESTY est composée d’une application mobile développée avec Flutter, destinée principalement aux visiteurs, utilisateurs et agents de scan. Elle comprend également des interfaces web développées avec React.js, notamment le portail partenaire et le tableau de bord administrateur.


Le Back-End FESTY, développé avec Spring Boot, représente le service applicatif central. Il expose des API REST sécurisées consommées par l’application mobile et les interfaces web. Il assure le traitement métier, la gestion des autorisations, l’accès aux données et l’intégration avec les services externes.


La base de données PostgreSQL [21] est hébergée dans l’environnement Supabase. Elle permet de stocker les données principales de la plateforme, telles que les utilisateurs, les partenaires, les événements, les billets, les paiements et les opérations de scan. L’extension PostGIS [23] est utilisée pour gérer les données géographiques liées à la localisation des événements.


La plateforme intègre aussi plusieurs services externes. Stripe assure les paiements, les remboursements et les opérations financières. Firebase Cloud Messaging [26] prend en charge les notifications push. Brevo [27] est utilisé pour les SMS de vérification OTP ainsi que pour les emails transactionnels, et Bunny.net [28] pour la gestion ou la diffusion des fichiers et médias.


La figure suivante présente l’architecture physique de la plateforme FESTY.


\begin{figure}[H]
\centering
\includegraphics[width=0.95\textwidth]{\detokenize{images/chapitre2/Figure 2.3 Architecture physique de la plateforme FESTY.png}}
\caption{Architecture physique de la plateforme FESTY.}
\label{figure:2-3}
\end{figure}


Cette architecture montre que le Back-End joue un rôle central dans la coordination du système. Il assure les échanges entre les interfaces clientes, la base de données et les services externes, tout en appliquant les règles métier et les contrôles de sécurité nécessaires.


\subsection{Architecture logique}


L’architecture logique de FESTY repose sur un Back-End unique développé avec Spring Boot et organisé sous forme de monolithe modulaire. Ce choix permet de conserver une solution simple à développer et à déployer, tout en structurant le code en modules métier cohérents.


Afin d’assurer cette organisation, nous avons utilisé Spring Modulith. Cette approche permet de découper le Back-End selon les principaux domaines fonctionnels de la plateforme, tels que l’authentification, les utilisateurs, les partenaires, les événements, la billetterie, le paiement, le scan, les notifications, les interactions, les stars et recommandations ainsi que l’administration.


Chaque module possède un périmètre fonctionnel précis et expose uniquement les services nécessaires aux autres modules. Cette séparation permet de réduire le couplage entre les composants, d’améliorer la lisibilité du code et de faciliter la maintenance de la plateforme.


Les échanges entre modules peuvent être réalisés à travers des interfaces ou des événements applicatifs. Par exemple, lorsqu’une action importante est effectuée dans un module, un événement peut être publié afin d’être traité par un autre module sans créer une dépendance forte entre les deux. Cette organisation permet de rapprocher certains avantages des microservices, tout en conservant la simplicité de déploiement d’un monolithe.


Les principaux modules logiques de FESTY sont présentés dans le tableau suivant.


\begin{table}[H]
\centering
\caption{Modules logiques de la plateforme FESTY}
\label{tableau:2-14}
\renewcommand{\arraystretch}{1.2}
\begin{tabularx}{\textwidth}{|>{\raggedright\arraybackslash}X|>{\raggedright\arraybackslash}X|}
\hline
\rowcolor[HTML]{ECF4FF}
\textbf{Module} & \textbf{Rôle principal} \\
\hline
Authentification & Gestion de l’inscription, de l’authentification, des jetons d’accès et de la sécurisation des comptes. \\
\hline
Utilisateurs & Gestion des profils, préférences et centres d’intérêts. \\
\hline
Partenaires & Gestion des candidatures, profils professionnels et documents de vérification. \\
\hline
Événements & Gestion de la création, consultation, recherche et cycle de vie des événements. \\
\hline
Billetterie & Gestion des types de billets, des réservations et des tickets numériques. \\
\hline
Paiement & Gestion des paiements, remboursements, soldes et opérations financières. \\
\hline
Scan & Gestion du contrôle d’accès, des sessions de scan et des validations d’entrée. \\
\hline
Notifications & Gestion des notifications et préférences de notification. \\
\hline
Interactions & Gestion des avis, favoris, chat et signalements. \\
\hline
Stars et recommandations & Gestion des artistes, intérêts utilisateurs et recommandations personnalisées. \\
\hline
Administration & Supervision de la plateforme, gestion du catalogue artistique, traitement des opérations financières et contrôle des opérations sensibles. \\
\hline
\end{tabularx}
\end{table}


La figure suivante présente une vue simplifiée de l’architecture logique modulaire du Back-End FESTY.


\begin{figure}[H]
\centering
\includegraphics[width=0.95\textwidth]{\detokenize{images/chapitre2/Figure 2.4 Architecture logique modulaire de FESTY.png}}
\caption{Architecture logique modulaire de FESTY.}
\label{figure:2-4}
\end{figure}


Cette architecture permet de mieux organiser les responsabilités internes de la plateforme. Elle facilite également l’évolution future du système, car chaque module peut être modifié ou enrichi avec un impact limité sur les autres parties.


\section{Conclusion}


Dans ce chapitre, nous avons présenté la phase d’analyse et de préparation du projet FESTY. Nous avons identifié les acteurs principaux de la plateforme, puis défini les besoins fonctionnels et non fonctionnels du système. Nous avons également présenté le backlog produit et la planification du projet en quatre sprints.


Nous avons ensuite introduit les deux diagrammes de cas d’utilisation globaux, ainsi que l’environnement de travail, les technologies utilisées et l’architecture de la plateforme. Le chapitre suivant sera consacré à l’étude et à la réalisation du premier sprint, portant sur l’authentification et la gestion des accès.


